\section{Computational methods}
\label{sec:methods}

\subsection{Configuration generation and selection}
\label{sec:configurations}

The reference database was constructed in two stages: exploration of Ti--Al configurations with a pretrained potential, followed by DFT labeling of a compact representative subset.
The medium MACE-MPA-0 foundation model was used for the exploration stage,\cite{Batatia2025} with molecular dynamics and structural optimization implemented through the Atomic Simulation Environment (ASE).\cite{Larsen2017}
The configuration pool combined liquid trajectories, configurations sampled during relaxation from liquid snapshots, and trajectories initiated from substitutional solid solutions.
The crystalline parent structures included hcp and bcc Ti, $\gamma$-TiAl, and $\alpha_2$-Ti$_3$Al.
Off-stoichiometric configurations were generated by randomly replacing atomic species on the parent lattices.
These structures were random substitutional realizations; no optimization of special quasirandom structures was performed.

The solid-solution sampling runs were carried out at 1000~K for 10~ps, whereas the liquid-generation scripts used 2000~K and included alternating MD and relaxation stages.
The exploration dynamics used a 1~fs time step and a target pressure of 1~bar.
Relaxation trajectories supplemented the finite-temperature configurations with lower-energy local environments.
MACE energies and forces were used to generate this pool; the targets for fine-tuning were subsequently obtained from DFT.
% The initially described 1 ps interval between liquid relaxation starts is
% not asserted here: archived script variants use 100 or 500 MD steps.

Representative structures were selected with ASAP\cite{Cheng2020} using smooth overlap of atomic positions (SOAP) descriptors.\cite{Bartok2013}
The SOAP expansion used a radial cutoff of 3.2~\AA, a Gaussian width of 0.4~\AA, four radial basis functions, and angular terms up to $l=3$, with Ti--Al cross-species contributions included.
Atomic descriptors were averaged over each periodic configuration.
Farthest-point sampling in this descriptor space reduced the pool to 1000 configurations.
The selected set contained 880 structures with 216 atoms, 104 with 250 atoms, and 16 with 256 atoms, and included both pure elements and alloy compositions.
% FPS is inferred from the archived ASAP command, which leaves the algorithm
% at the documented default. Record the installed ASAP version before release.

\subsection{Electronic-structure reference calculations}
\label{sec:dft}

Single-point DFT calculations were performed at the selected geometries using VASP\cite{Kresse1996} and the projector augmented-wave method.\cite{Blochl1994,Kresse1999}
Three reference descriptions were considered: PBE, PBEsol, and \rscan.
The supplied PAW datasets were Ti\_pv, with ten valence electrons, and Al, with three valence electrons.
All three reference datasets were generated with a common numerical protocol in VASP~6.3.2 using non-spin-polarized calculations, a plane-wave cutoff of 550~eV, and an electronic convergence threshold of $10^{-6}$~eV.
First-order Methfessel--Paxton smearing with a width of 0.1~eV was employed.
The calculations used the \texttt{Accurate} precision setting, automatic real-space projection, and nonspherical contributions to the PAW exchange-correlation terms.
Explicit configuration-dependent reciprocal-space grids contained 8--12 full-grid points, reduced to 4--6 points in the supplied \texttt{KPOINTS} files.
These grids excluded $\Gamma$ and were identical for a given configuration across all three functionals.
The ions and simulation cells were held fixed during labeling.

The resulting datasets supplied total energies, atomic forces, and virials for fine-tuning.
Their sizes and training/validation partitions are listed in Table~\ref{tab:datasets}.
All retained structures could be matched to the common set of 1000 selected configurations; 992 configurations were present in all three labeled datasets.
There was no configuration overlap between the training and validation subsets of an individual functional.
The partitions were constructed separately, however, and were not identical across functionals.
Accordingly, errors on these validation subsets characterize approximation of the respective DFT datasets and are not treated as errors on a shared independent test set.

\begin{table}[tb]
\caption{Numbers of configurations in the fine-tuning datasets. The selected pool contains 1000 configurations.}
\label{tab:datasets}
\begin{ruledtabular}
\begin{tabular}{lrrr}
Reference & Total & Training & Validation \\
\hline
PBE      & 999 & 850 & 149 \\
PBEsol   & 998 & 849 & 149 \\
\rscan   & 993 & 845 & 148 \\
\end{tabular}
\end{ruledtabular}
\end{table}

\subsection{Fine-tuning DPA4-Neo}
\label{sec:finetuning}

The three Ti--Al potentials were initialized from the pretrained DPA4-Neo model,\cite{Li2026DPA4} using the checkpoint \texttt{DPA4-Neo-OMat24-v20260805.pt}.
Training used the PyTorch implementation in DeePMD-kit.
All three fine-tuning runs used the same hyperparameters, with a 6.0~\AA{} cutoff, two interaction blocks, 32 channels, 16 radial functions, and maximum angular order $l_{\max}=3$.
Energy-derived forces preserve the conservative form of the potential during adaptation.

Low-rank adaptation (LoRA)\cite{Hu2021} was enabled with rank 16 and scaling parameter 16.
The training schedule comprised 50 epochs, with a cosine learning-rate decay from $10^{-4}$ to $10^{-6}$.
The warm-up occupied 0.3\% of the schedule and started at 0.2 times the initial learning rate.
Optimization used HybridMuon with a weight decay of 0.01 and gradient-norm clipping at 5.
The loss combined mean-absolute-error terms for energy, forces, and virial, with constant prefactors of 20, 20, and 5, respectively; the force-norm option was enabled.
Automatic batch sizing was set to \texttt{auto:128}.
The model parameters used single precision, with automatic mixed precision enabled for training.
Source-checkpoint and frozen-export identities are recorded by cryptographic hashes for the held-out evaluation.
% Do not interpret "auto:128" as 128 structures per batch.
% Exact loss normalization / force-norm implementation and optimizer version
% must be pinned to the installed DPA4 code before reproducing the training.

The checkpoint retained for evaluation was selected separately from the duration of the training schedule.
For PBE, a widening training--validation gap in the logged force errors raised a concern about overfitting and motivated comparison of earlier and later checkpoints.
The checkpoint at approximately 5000 optimization steps was retained because it gave better agreement with an independently calculated PBE DFT phonon density of states (DOS) for hcp Ti than the final checkpoint at approximately 23\,400 steps.
The selection criterion was the DFT DOS, rather than the experimental phonon dispersion points.
The retrospective analysis in the Supporting Materials confirms this preference over the 20\,000- and 23\,400-step checkpoints when all three models are evaluated at the exact DFT reference geometry, without relaxation.
The same analysis shows that validation and held-out energy and force errors decrease overall beyond 5000 steps; the growing gap in the learning curves therefore does not establish a general deterioration of predictive accuracy.
For PBEsol and \rscan{}, the final checkpoints were retained after the learning curves reached a plateau and the final models gave the preferred property description.
Thus, the PBE model represents a deliberate choice based on phonon performance, rather than selection of the lowest energy error or use of the same optimization step for all three functionals.
Agreement with the Ti DFT DOS consequently forms part of model selection and is not a wholly independent assessment of the selected PBE model.
The common held-out DFT benchmark described below was not used to select these checkpoints.
Learning curves, the checkpoint-dependent errors, composition and volume diagnostics, and phonon calculations at common and separately optimized geometries are documented in the Supporting Materials.

\subsection{Common held-out DFT benchmark}
\label{sec:testset}

A fixed 50-configuration holdout was selected independently of model errors: 25 configurations from a separate 200-structure liquid archive and 25 from trajectories initialized as hcp, bcc, and $\gamma$-TiAl solid solutions (8, 9, and 8 configurations, respectively).
Selection used composition-balanced farthest-point sampling of normalized, configuration-averaged SOAP descriptors within each family, with a 5~\AA{} cutoff, Gaussian width 0.5~\AA, four radial functions, and $l_{\max}=3$.
No selected geometry matched the supplied training or validation geometries under species-preserving periodic matching, atom permutations, and common translations, using a $5\times10^{-5}$~\AA{} tolerance and fixed cell axes.
An additional near-duplicate screen used a 0.1~\AA{} RMSD threshold together with less than 1\% cell deformation.
This establishes a geometrical holdout from the fine-tuning data; independence of the underlying liquid trajectories and absence of overlap with foundation-model pretraining are not established.
The liquid designation records the source pool, rather than a single equilibrium temperature.

The test configurations were labeled without relaxation using VASP~6.6.1, a 550~eV cutoff, $10^{-6}$~eV electronic convergence, and first-order Methfessel--Paxton smearing of width 0.1~eV.
The PAW datasets were unchanged; reciprocal-space grids were generated with kpLib~1.1.1 using a minimum-distance criterion of 35~\AA.
For each geometry, all three functionals used identical PAW files and reciprocal-space grids.
These test calculations retain the principal numerical settings of the labeling protocol, while the VASP and grid-generator versions differ from those of the original datasets.
Only configurations with converged labels for all three functionals enter a comparison, with geometry and file identities checked before analysis.
Frozen exports of the supplied fine-tuned checkpoints were evaluated on the same geometries; the unchanged pretrained checkpoint was evaluated once and its predictions compared separately with each functional.
No model parameters or energy offsets were fitted to the holdout.

Energy errors compare raw model energies with VASP free energies (\texttt{TOTEN}), divided by the number of atoms.
Force MAE and RMSE pool all Cartesian components; stress errors pool the six independent tensor entries in GPa, with model stress obtained as minus virial divided by volume.
Descriptive 95\% percentile intervals use 2000 bootstrap resamples of whole configurations, retaining the force components within each configuration together and using matched draws across functionals.
These intervals are conditional on the currently completed subset and do not account for unknown trajectory correlations or repeated inspection as the test set grows.

\subsection{Thermal density and liquid mixing thermodynamics}
\label{sec:thermodynamics}

Density and mixing thermodynamics were evaluated using periodic NPT simulations in ASE with the fine-tuned potentials.
The integration time step was 1~fs and the external pressure was 1~bar.
The stored MD implementation used the Melchionna NPT integrator, a thermostat time parameter of 25~fs, and isotropic cell fluctuations.
The barostat parameter was set to $B\tau^2$, where $\tau=25$~fs and $B$ was the bulk-modulus estimate supplied by the workflow.
For liquid Ti--Al mixtures, this estimate was interpolated between 110~GPa for Ti and 76~GPa for Al.
% The barostat interpolation values are checked against utils/bulk_modulus.py.

For the temperature-dependent density of Ti, 250-atom cells were initialized from $\alpha$-Ti at 300, 500, 700, 900, and 1100~K and from $\beta$-Ti at 1200, 1400, 1600, 1800, and 1900~K.
The stored solid-state runs comprised 20~ps of equilibration and 30~ps of production.
Liquid Ti was simulated in 200-atom cells between 2000 and 2500~K in 100~K increments.
Liquid preparation involved 10~ps at $1.5T$, followed by 30~ps at the target temperature and a further 30~ps of production.
Production configurations were saved every 0.1~ps.
The initial crystalline phase labels specify the preparation procedure; interpretation of a high-temperature density point also requires checking whether that phase persisted during the trajectory.

For the alloy liquids, we define the composition as the Al atomic fraction,
\begin{equation}
 x = x_{\mathrm{Al}} = \frac{N_{\mathrm{Al}}}{N_{\mathrm{Al}}+N_{\mathrm{Ti}}}.
\end{equation}
Calculations used 200 atoms at $x=0,0.1,\ldots,1$ and temperatures of 1873, 2000, and 2500~K, following the same liquid preparation and production protocol.
The 1873 and 2000~K series provide the principal comparison; the 2500~K series extends the temperature range.
Density and mixing enthalpy were obtained from the same production trajectories.
The mass density was calculated as the time average $\rho=\langle M/V\rangle$, where $M$ is the total mass of the simulated cell.

For each model, the molar enthalpy was evaluated as
\begin{equation}
 H_m(T,p,x)=\frac{N_{\mathrm A}}{N}
 \left[\langle E_{\mathrm{pot}}+E_{\mathrm{kin}}\rangle+
 p_{\mathrm{ext}}\langle V\rangle\right],
 \label{eq:enthalpy}
\end{equation}
where $N_{\mathrm A}$ is Avogadro's constant and $N$ is the number of atoms.
The mixing enthalpy is
\begin{align}
 \Delta H_{\mathrm{mix}}(T,p,x)={}&H_m(T,p,x)\nonumber\\
 &-xH_m^{\mathrm{Al}}(T,p)-(1-x)H_m^{\mathrm{Ti}}(T,p).
 \label{eq:mixing}
\end{align}
The pure-liquid reference values were evaluated with the same potential at the same temperature and pressure.
Equation~\eqref{eq:enthalpy} uses the prescribed external pressure; the ASE stress convention is $p=-\mathrm{Tr}\,\boldsymbol{\sigma}/3$.
The enthalpies were reprocessed from the stored energies and volumes using this definition.
At a common temperature the classical kinetic contribution cancels in the thermodynamic mixing limit.
The excess molar volume was calculated analogously from $V_m=N_{\mathrm A}\langle V\rangle/N$:
\begin{equation}
 V_m^{\mathrm{ex}}(x)=V_m(x)-xV_m^{\mathrm{Al}}-(1-x)V_m^{\mathrm{Ti}}.
\end{equation}
Liquid-state values below the corresponding liquidus describe undercooled liquids and must be distinguished from stable equilibrium phase mixtures.

Statistical variability was assessed by subdividing production trajectories into contiguous time blocks.
The preliminary thermodynamic analysis used six 5~ps blocks, with additional block lengths examined as a correlation-time diagnostic.
Uncertainty in a mixing quantity includes the uncertainties of both pure-component reference values; reuse of these references also correlates errors between compositions.
The final uncertainty estimates require stabilization with respect to block length and sampling time.

\subsection{Harmonic lattice dynamics}
\label{sec:phonons}

Harmonic phonons of $\alpha$-Ti and $\alpha_2$-Ti$_3$Al were calculated using the finite-displacement implementation in ASE.
Before evaluating force constants, the structures were optimized with BFGS, preserving crystal symmetry and allowing hydrostatic cell strain through a cell filter.
The convergence target was $10^{-4}$~eV\,\AA$^{-1}$.
The actual production calculations used $5\times5\times5$ supercells: 250 atoms for $\alpha$-Ti and 1000 atoms for Ti$_3$Al.
Positive and negative Cartesian displacements of the atoms in the reference cell were used to construct force constants, followed by enforcement of the acoustic sum rule.
The archived production commands specify a displacement amplitude of 0.01~\AA{}, a $50\times50\times50$ phonon mesh for the DOS, and a Gaussian width of $2\times10^{-4}$~eV (approximately 0.048~THz).
Phonon dispersions and densities of states were obtained from the corresponding dynamical matrices.
These are harmonic calculations about the optimized reference structures and do not include finite-temperature anharmonic renormalization.

An independently calculated PBE phonon reference for hcp Ti was also available.
It used Phonopy\cite{Togo2015} with a $4\times4\times3$ supercell containing 96 atoms and a displacement amplitude of 0.01~\AA.
The static VASP force calculations used a 500~eV cutoff, a $4\times4\times4$ $\Gamma$-centered electronic mesh, a $10^{-8}$~eV electronic convergence threshold, and a smearing width of 0.2~eV.
Its density of states was sampled on a $40\times40\times30$ phonon mesh and broadened with a Gaussian width of 0.05~THz.
For comparisons of spectral weight, densities of states are normalized per atom so that their frequency integral is three.

\draftnote{The separate DFT Ti$_3$Al dataset has an unresolved geometry assignment (the supplied audit identifies Cmcm symmetry), so it is not yet adopted as a verified reference for ideal D0$_{19}$ Ti$_3$Al.}

\subsection{Solid--liquid coexistence and melting of Ti}
\label{sec:melting}

The melting temperature of Ti was estimated by direct solid--liquid coexistence simulations in LAMMPS.\cite{Thompson2022}
A periodic 1296-atom cell, approximately $20\times20\times59$~\AA$^3$ at the beginning of production, contained a bcc solid slab and a liquid slab separated along $z$.
Periodic boundaries produced two interfaces.
The preparation included 5 or 10~ps of initial equilibration, 20~ps of heating one half of the sample to $1.5T$, and 10~ps of cooling that region to the target temperature while retaining a crystalline seed.
The complete system was then propagated for 50~ps in the NPT ensemble at zero external pressure.
The time step was 1~fs; thermostat and barostat damping times were 0.1 and 1~ps, respectively.
The three box lengths were allowed to fluctuate independently.
Configurations were saved every 0.5~ps.

The direction of interface motion was determined from spatially resolved structural order and the evolution of the solid-like fraction.
For a quantitative check, neighbor-averaged bond-order parameters $\bar q_6$ were calculated following Lechner and Dellago,\cite{Lechner2008} using the 14 nearest neighbors.
A threshold of $\bar q_6>0.25$ was used to identify solid-like atoms, with alternative thresholds and a bond-correlation criterion used to assess classification sensitivity.
Profiles were evaluated along the reduced coordinate $z/L_z$ to separate interface motion from changes in box length.
Growth and recession of the crystalline layer were used to bracket the melting temperature, with complete crystallization or melting providing additional endpoint checks.
The resulting temperature intervals describe the observed behavior of this finite cell over 50~ps; their half-widths are not interpreted as full statistical uncertainties.

\subsection{Liquid viscosity}
\label{sec:viscosity}

Shear viscosity was calculated from equilibrium stress fluctuations using the Green--Kubo relation.
The viscosity simulations used 512-atom periodic cells and the PBE- and \rscan{}-based potentials.
At 1873~K, the principal series covered Ti fractions
$c_{\mathrm{Ti}}=0.1,0.2,0.3,0.4,0.5,0.55,0.6,0.7,0.8,0.9$, and $1$.
Here $c_{\mathrm{Ti}}=1-x$; pure Al was not included in this STACIE series.
Additional pure-Ti calculations were performed at 1941, 2000, and 2100~K.
These are nominal alloy compositions, with integer atom counts used in each finite cell.

Starting snapshots were drawn from a 5~ps NPT preparation trajectory at the target temperature and 1~bar.
Each selected snapshot was assigned velocities and propagated for 1~ps in NVT before the viscosity production run.
Production trajectories were 20~ps long, with a 1~fs integration step and pressure-tensor recording at every step.
The NVT thermostat damping time was 0.2~ps.
Most state points used 15 trajectories.
At 1873~K, 50 trajectories were used for both $c_{\mathrm{Ti}}=0.55$ cases, for PBE at $c_{\mathrm{Ti}}=0.70$, and for \rscan{} at $c_{\mathrm{Ti}}=0.40$; the \rscan{} pure-Ti point used 30.
The PBE pure-Ti point at 2000~K also used 50 trajectories.
Each production trajectory retained the volume of its own starting snapshot.
The preparation procedure therefore requires attention to both equilibration and correlations between starting configurations when assessing statistical independence.

For an isotropic liquid, define five equivalent shear channels,
\begin{align}
 Q_1&=P_{xy},\quad Q_2=P_{xz},\quad Q_3=P_{yz},\nonumber\\
 Q_4&=\tfrac12(P_{yy}-P_{zz}),\nonumber\\
 Q_5&=\frac{P_{xx}-\tfrac12(P_{yy}+P_{zz})}{\sqrt{3}}.
\end{align}
The viscosity of a fixed-volume trajectory can then be expressed as
\begin{equation}
 \eta=\frac{V}{5k_{\mathrm B}T}\sum_{a=1}^{5}
 \int_0^\infty
 \left\langle\delta Q_a(0)\delta Q_a(t)\right\rangle\,dt,
 \label{eq:viscosity}
\end{equation}
where $\delta Q_a=Q_a-\langle Q_a\rangle$ and the pressure tensor includes both kinetic and configurational contributions.

The primary estimates were obtained with STACIE (STable AutoCorrelation Integral Estimator).\cite{Toraman2025}
The five channels from each trajectory were combined using its own prefactor $V/(k_{\mathrm B}T)$.
Pressure and time were converted to SI units before spectral analysis.
The low-frequency spectrum was modeled using the STACIE \texttt{LorentzModel}; the viscosity follows from its zero-frequency limit.
Original estimates were reproduced with STACIE~1.3.0; expanded-ensemble estimates use the supplied outputs.
Reported estimator standard deviations describe uncertainty within that spectral estimation procedure and do not include all finite-size, thermostat, preparation, or MLIP errors.

A time-decomposition analysis based on the approach of Zhang \textit{et al.}\cite{Zhang2015} was used for comparison.
This analysis averaged the three off-diagonal stress autocorrelation functions and fitted the running Green--Kubo integral with a weighted two-exponential form.
Comparisons between the two estimators used matching state points and trajectory sets.
Effective spectral sample size and sensitivity to trajectory subdivisions were examined as convergence diagnostics.
The additional 1773~K series currently has only time-decomposition estimates and is kept separate from the principal STACIE dataset.

\draftnote{The viscosity protocol describes the presently supplied trajectories. The agreed checks of sampling duration and preparation, and any additional STACIE processing at 1773~K, must be reflected here when those calculations are finalized.}
